\section{Metodyka}

\subsection{Scenariusze}

Zgodnie z~koncepcją mamy 24 scenariusze CRUD, po 6 na każdą operację.
Dodaliśmy do nich 3 scenariusze na danych JSON (tab.~\ref{tab:scen}).
Parametry każdego wykonania (identyfikatory, miasta, przedziały dat) wynikają
wyłącznie z~(zbioru, scenariusza, numeru iteracji), więc wszystkie bazy
wykonują dokładnie te same operacje na tych samych wierszach.

\begin{table}[H]
\centering
\small
\caption{Scenariusze. ,,F'' oznacza scenariusz z~fiksturą: dane, na których
działa operacja, są wstawiane przed pomiarem i~nie wliczają się do czasu.}
\label{tab:scen}
\begin{tabular}{lp{4.6cm}p{8.6cm}}
\toprule
ID & nazwa & co dokładnie jest mierzone \\
\midrule
C1 & Pojedyncza rezerwacja & transakcja: nowa rezerwacja + zajęcie 1 wolnego biletu \\
C2 & Rezerwacja z~kilkoma biletami & jak C1, 4 bilety \\
C3 (F) & Masowy import biletów & \num{10000} nowych biletów (pula sprzedażowa nowych wydarzeń) \\
C4 & Rejestracja klienta & wstawienie klienta z~preferencjami JSON \\
C5 & Wydarzenie z~cennikiem & wydarzenie i~cena dla każdego sektora obiektu \\
C6 (F) & Zapis płatności & płatność + zmiana statusu rezerwacji na opłaconą \\
\midrule
R1 & Rezerwacja po identyfikatorze & rezerwacja z~biletami, miejscami i~sektorami \\
R2 & Wolne miejsca na wydarzenie & wolne bilety wydarzenia z~sektorem, rzędem i~numerem, posortowane \\
R3 & Wydarzenia po dacie i~mieście & wydarzenia w~mieście w~30-dniowym oknie \\
R4 & Historia rezerwacji klienta & 50 ostatnich rezerwacji klienta z~liczbą biletów \\
R5 & Sprzedaż per wydarzenie & liczba i~suma cen sprzedanych biletów dla wydarzeń organizatora \\
R6 & 10 najpopularniejszych wydarzeń & agregacja po wszystkich sprzedanych biletach \\
\midrule
U1 & Zmiana statusu rezerwacji & rezerwacja $\to$ opłacona \\
U2 & Zajęcie miejsca (blokada) & warunkowe \texttt{wolny}$\to$\texttt{zablokowany} z~terminem blokady \\
U3 & Masowa zmiana cen w~sektorze & cena $\times 1{,}1$ dla wolnych biletów sektora + cennik \\
U4 & Zmiana danych klienta & e-mail, telefon, scalenie pól w~preferencjach JSON \\
U5 & Zmiana terminu wydarzenia & przesunięcie daty o~tydzień \\
U6 (F) & Wygaszenie nieopłaconych rezerwacji & \num{1000} przeterminowanych rezerwacji $\to$ wygasłe, zwolnienie ich biletów \\
\midrule
D1 & Anulowanie rezerwacji & usunięcie rezerwacji i~płatności, zwolnienie biletów \\
D2 & Rezerwacja wraz z~biletami & usunięcie rezerwacji, płatności i~jej biletów \\
D3 (F) & Masowe usunięcie wygasłych & usunięcie \num{1000} wygasłych rezerwacji \\
D4 & Usunięcie konta klienta & klient, jego rezerwacje i~płatności, zwolnienie biletów \\
D5 (F) & Usunięcie odwołanego wydarzenia & wydarzenie, jego bilety (całe obiekty) i~cennik \\
D6 (F) & Usunięcie danych starszych niż data & wydarzenia sprzed 2 lat z~biletami, cennikami, rezerwacjami i~płatnościami \\
\midrule
J1 & Filtr po polu JSON & wydarzenia o~danym gatunku z~danym tagiem w~metadanych \\
J2 & Agregacja po polu JSON & liczba i~suma rezerwacji z~7 dni pogrupowana po \texttt{metadane.kanal} \\
J3 & Modyfikacja tablicy w~JSON & dopisanie tagu do \texttt{metadane.tagi} \\
\bottomrule
\end{tabular}
\end{table}

\subsection{Przebieg pomiaru}

Dla każdego rozmiaru zbioru i~każdej bazy plan wygląda tak:
\begin{enumerate}
  \item reset schematu, import danych, pomiar rozmiaru;
  \item wszystkie scenariusze bez indeksów drugorzędnych: 1 przebieg rozgrzewkowy
  i~5 mierzonych; plany zapytań;
  \item utworzenie indeksów (czas tworzenia jest osobnym wynikiem), pomiar rozmiaru;
  \item wszystkie scenariusze z~indeksami, plany zapytań;
  \item test współbieżności dla każdej strategii rezerwacji;
  \item reset (zwolnienie miejsca).
\end{enumerate}

Czas jest mierzony w~serwisie bazy, wokół samej operacji na bazie
(\texttt{time.Now()} przed i~po wywołaniu \texttt{Engine.Run}). Fikstura
(np. \num{1000} rezerwacji do usunięcia w~D3) jest wstawiana przed
rozpoczęciem pomiaru. Runner mierzy niezależnie czas całego wywołania gRPC,
dzięki czemu narzut sieci i~serializacji jest znany
(rozdz.~\ref{sec:grpc}), ale nie wpływa na porównanie baz. Wynik scenariusza
to średnia z~5 powtórzeń. W~plikach z~danymi są też mediana, odchylenie
standardowe, p95, minimum i~maksimum. Jedno wykonanie ma limit 60~s. Po
przekroczeniu limitu dalsze powtórzenia są pomijane, a~w~tabeli pojawia się
,,t/o''.

Kolejne powtórzenia działają na różnych wierszach (inny numer iteracji),
a~przebieg z~indeksami używa innego zakresu iteracji niż przebieg bez
indeksów. Scenariusze rezerwujące miejsca nie trafiają więc na bilety zajęte
we wcześniejszym przebiegu. Plany zapytań powstają przez
\texttt{EXPLAIN (ANALYZE, BUFFERS)} w~PostgreSQL, \texttt{ANALYZE FORMAT=JSON}
w~MariaDB (oba w~transakcji wycofywanej po pomiarze) oraz
\texttt{explain} z~\texttt{executionStats} w~MongoDB. Dla Valkey zapisujemy
listę wykonywanych poleceń.

\subsection{Test współbieżności}

Test odpowiada na pytanie z~hipotezy H1. W~serwisie bazy startuje
$N \in \{1, 10, 100, 1000\}$ gorutyn, każda z~własną sesją bazy. Wszystkie
próbują zarezerwować te same 500 biletów jednego wydarzenia, każda w~innej
losowej kolejności. Runda kończy się, gdy wszystkie bilety są zajęte. Mierzymy
czas rundy, przepustowość (udane rezerwacje na sekundę), liczbę konfliktów
(próba trafiła na zajęte miejsce albo transakcja została przerwana)
i~percentyle czasu pojedynczej próby. Po każdej rundzie serwis sprawdza
w~bazie, że istnieje dokładnie 500 rezerwacji i~500 zajętych biletów, czyli
że żadne miejsce nie zostało sprzedane dwa razy. Każdą konfigurację powtarzamy
3 razy.

Strategie rezerwacji:
\begin{itemize}
  \item PostgreSQL i~MariaDB: \texttt{cond} (warunkowy \texttt{UPDATE ... WHERE status='wolny'}),
  \texttt{forupdate} (\texttt{SELECT ... FOR UPDATE}, sprawdzenie, \texttt{UPDATE}),
  \texttt{nowait} (jak wyżej, ale zablokowany wiersz od razu oznacza konflikt),
  \texttt{serializable} (izolacja SERIALIZABLE i~ponawianie po błędzie serializacji lub zakleszczeniu);
  \item MongoDB: \texttt{atomic} (warunkowy \texttt{updateOne} na dokumencie biletu, potem
  \texttt{insertOne} rezerwacji) i~\texttt{txn} (obie operacje w~transakcji, \texttt{WriteConflict} to konflikt);
  \item Valkey: \texttt{lua} (sprawdzenie i~zapis w~jednym skrypcie),
  \texttt{setnx} (klucz blokady \texttt{SET NX PX}, czyli ,,blokada miejsca z~TTL''
  z~koncepcji, potem \texttt{MULTI}) i~\texttt{watch} (optymistyczne
  \texttt{WATCH}/\texttt{MULTI}/\texttt{EXEC} na zbiorze wolnych miejsc).
\end{itemize}
Pierwsza strategia z~listy jest domyślna i~to ona jest używana do weryfikacji H1.

\subsection{Środowisko}

\begin{itemize}
  \item Komputer: Apple M5 (10 rdzeni), 24\,GB RAM, macOS 27, Docker Desktop
  (maszyna wirtualna Linux 6.12, 10 vCPU, 8\,GB RAM).
  \item Każda z~trzech baz dyskowych dostaje 1\,GB pamięci podręcznej
  (\texttt{shared\_buffers}, \texttt{innodb\_buffer\_pool\_size},
  \texttt{wiredTigerCacheSizeGB}). Valkey trzyma cały zbiór w~RAM, bez RDB i~AOF.
  \item Trwałość zostaje domyślna: \texttt{fsync} i~\texttt{synchronous\_commit}
  w~PostgreSQL, \texttt{innodb\_flush\_log\_at\_trx\_commit=1} w~MariaDB,
  \emph{write concern} \texttt{majority} z~dziennikiem w~MongoDB. Valkey
  z~definicji nie zapisuje danych na dysk, co trzeba brać pod uwagę przy
  porównaniu zapisów.
  \item Poziomy izolacji są domyślne dla silnika: READ COMMITTED w~PostgreSQL,
  REPEATABLE READ w~MariaDB, \emph{snapshot} w~transakcjach MongoDB.
  \item Pule połączeń do 1100 sesji, \texttt{max\_connections=1200}.
  \item Przy zbiorach 1 i~10 mln biletów, na początku każdego segmentu,
  restartowane są kontenery pozostałych trzech baz. Mierzona baza ma wtedy do
  dyspozycji całą pamięć maszyny wirtualnej, a~pamięć podręczna innych silników
  jej nie zabiera.
\end{itemize}
